% TOP_PROBLEM: {"id": 2645, "title": "Kourovka Notebook Problem 21.136", "category_id": 17, "category": {"id": 17, "name": "group_theory", "display_name": "Group Theory", "description": "Problems about groups, group actions, representations, and related algebraic structures.", "slug": "group-theory", "order_index": 17, "created_at": "2026-07-31T15:26:25.671Z"}, "status": "open", "problem_number": "KOU-21.136", "set_id": 10}
% FIRST_POSTED: 2026-10-01
% ENTRY_KIND: statement_only
% UnsolvedMath ID 2645; source catalogue code KOU-21.136
% !TeX program = lualatex
% Definition and sources only; this file is not an LLM solution attempt.
\documentclass[11pt,a4paper]{article}
\usepackage[margin=25mm]{geometry}
\usepackage{fontspec}
\setmainfont{lmroman10-regular.otf}[BoldFont=lmroman10-bold.otf,
 ItalicFont=lmroman10-italic.otf,BoldItalicFont=lmroman10-bolditalic.otf]
\usepackage{amsmath,amssymb,mathtools}
\usepackage{unicode-math}
\setmathfont{latinmodern-math.otf}
\AtBeginDocument{\let\setminus\smallsetminus}
\usepackage{xurl,hyperref}
\hypersetup{colorlinks=true,urlcolor=blue}
\setcounter{secnumdepth}{0}
\setlength{\parindent}{0pt}
\setlength{\parskip}{5pt}
\setlength{\emergencystretch}{3em}
\begin{document}
\section*{Kourovka Notebook Problem 21.136}\label{2645}
{\small UnsolvedMath ID 2645 · KOU-21.136 · Group Theory · Kourovka Notebook - New Problems, Issue 21}
\subsection{Definitions and mathematical statement}
A profinite group is a compact Hausdorff topological group that is an inverse limit of finite discrete groups. Conjugacy in $G$ is the relation $x\sim y$ when $y=g^{-1}xg$ for some $g\in G$; the conjugacy class of $x$ is $x^G=\{g^{-1}xg:g\in G\}$.

An element has infinite order if $x^n\ne1$ for every positive integer $n$. Put
\[
 \mathcal C_\infty(G)=\{x^G:x\in G\text{ has infinite order}\}.
\]
Since element order is constant on a conjugacy class, this is a well-defined set of classes. Let $\mathfrak c=2^{\aleph_0}$ denote the cardinality of the real numbers.

If $|\mathcal C_\infty(G)|<\mathfrak c$, must $G$ be torsion? Here torsion means that every individual element has some finite order, with no common exponent bound required. Equivalently, does the existence of even one infinite-order element force at least continuum many distinct conjugacy classes consisting of infinite-order elements?

The hypothesis counts classes, not elements, so a single class may itself be infinite or uncountable. The inequality is strict and is a comparison of cardinals, without an assumption about the continuum hypothesis. No finite-generation, metrizability, or countability condition is imposed on the profinite group. All conjugations and all element orders are taken in the given group $G$.
\subsection{Short English statement}
If a profinite group has any infinite-order element, must it have at least continuum many conjugacy classes of such elements?
\subsection{Sources}
\begin{itemize}
\item[S1] ULAM.AI and the UnsolvedMath Contributors, supplied problems.json, record 2645 (KOU-21.136), Kourovka Notebook - New Problems, Issue 21. Catalogue identity and source transcription; CC BY 4.0.
\url{https://huggingface.co/datasets/ulamai/UnsolvedMath}.
\item[S2] The Kourovka Notebook, 21st edition, version 47 (30 September 2026), new problems, Problem 21.136. Expanded formulation of the supplied catalogue statement.
\url{https://arxiv.org/pdf/1401.0300v47}.
\end{itemize}
\end{document}
